\addtocontents{toc}{\protect\Needspace{3\baselineskip}}
\originalchapter{8.821 课程详讲：热态、相变与化学势}
本章对应第 22 讲后半和第 23 讲. 平直边界只含温度尺度, 球形边界还含半径, 因而两者的热鞍点结构不同. 最后以明确定义的 Einstein--Maxwell 截断求一个有限密度解, 不把它视为所有 IIB 有荷相的完整分类. 单位 $\hbar=c=\kb=1$.

\originalsection{热态需要怎样的体背景}
边界正则配分函数为 $Z(\beta)=\operatorname{Tr}\ee^{-\beta H}$, 欧氏时间周期 $\beta$. 玻色场周期、费米场反周期, 它们是热边界条件, 不是零温超对称指数的边界条件. 候选体鞍点应具有同一渐近边界、相同欧氏周期, 并可按热力学比较壳上作用量.

平直边界上, 一种形式候选是把纯 Poincaré AdS 的欧氏时间直接粘成圆. 它的\textbf{局部曲率仍是 AdS 恒曲率}, 不能说 $z\to\infty$ 使其 Ricci 曲率本身发散. 但圆的固有周长 $\beta R/z$ 趋零, 在该非可收缩识别中会进入弦尺度, 缠绕模和相应热不稳定不能用仅含低能场的气体描述忽略. 与此相比, 黑膜欧氏时间圆在视界处像正则极坐标原点收缩, 几何是一个雪茄形, 拓扑和正则性不同.

对平直 $\mathbb R^{d-1}$ 边界, 保持平移和旋转的真空 Einstein 解可取
\[
\dd s^2=\frac{R^2}{z^2}
\left[-f(z)\dd t^2+\dd\mathbf x^2+\frac{\dd z^2}{f(z)}\right],
\qquad f(z)=1-\left(\frac z{z_h}\right)^d.
\]
空间对角 Einstein 分量给 $zf'-df=-d$, 其积分是常数1加 $z^d$ 项, 视界条件确定上述系数. 其余分量对这套假设形式也成立. $z_h$ 不是另一份独立的 CFT 尺度, 正则性将它与温度联系.

近视界的固有径向坐标可取 $\rho=2R\sqrt{(z_h-z)/(dz_h)}$, 欧氏径向时间部分成为
\[
\dd\rho^2+\left(\frac d{2z_h}\right)^2\rho^2\dd t_E^2.
\]
因此 $t_E$ 周期为 $4\pi z_h/d$, 温度 $T=d/(4\pi z_h)$. 对 $d=4$, $z_h=1/(\pi T)$; 温度越高, 体中视界越靠近边界紫外端.

\originalsection{面积给熵密度, 热力学给能量与压强}
坐标空间体积 $V=\int\dd^{d-1}x$ 的视界面积为 $A=(R/z_h)^{d-1}V$, 熵密度
\[
s=\frac1{4G_{d+1}}\left(\frac R{z_h}\right)^{d-1}.
\]
对 $\mathcal N=4$ 的 $d=4$ 字典 $G_5/R^3=\pi/(2N^2)$,
\[
s=\frac{\pi^2}2N^2T^3.
\]
取平直真空热自由能为零, $s=-\partial f_{\rm th}/\partial T$, 从而
\[
f_{\rm th}=-\frac{\pi^2}8N^2T^4,
\quad p=-f_{\rm th}=\frac{\pi^2}8N^2T^4,
\quad e=f_{\rm th}+Ts=\frac{3\pi^2}8N^2T^4.
\]
$e=3p$ 与平直 CFT 的无迹应力相容. 这些系数来自强耦合、大 $N$、领先 Einstein 近似; 不能把它们称为任意 $\lambda$ 的精确热方程.

自由 SYM 有8个玻色和8个费米在壳自由度, 所以
\[
s_{\rm free}=\frac{2\pi^2}{45}
\left(8+\frac78\,8\right)(N^2-1)T^3
\simeq\frac{2\pi^2}3N^2T^3.
\]
同一大 $N$ 取法下 $s_{\rm strong}/s_{\rm free}=3/4$. 它说明相互作用改变了热自由度的有效系数, 不是有四分之一基本场从 Lagrangian 消失. 其他 CFT 的比值需各自模型和字典, 不由这一次计数普适固定.

\originalsection{球形边界允许热 AdS 与大小黑洞}
边界空间是半径 $R_b$ 的球时, 物理可依赖无量纲 $R_bT$. 以下取自然共形代表 $R_b=R$. 欧氏全局 AdS 的时间圆固有周长为 $\beta\sqrt{1+r^2/R^2}$, 在中心的最小值是 $\beta$ 而非0. 当它仍大于有关弦尺度, 低温热 AdS 气体是合适鞍点.

球形黑洞则取
\[
\dd s^2=-f(r)\dd t^2+\frac{\dd r^2}{f(r)}+r^2\dd\Omega_{d-1}^2,
\quad f(r)=1+\frac{r^2}{R^2}-\frac\mu{r^{d-2}}.
\]
由 $f(r_h)=0$ 得 $\mu=r_h^{d-2}(1+r_h^2/R^2)$, 由正则性得
\[
T=\frac{dr_h^2+(d-2)R^2}{4\pi R^2r_h}.
\]
对 $d>2$, 最低温度出现在 $r_h/R=\sqrt{(d-2)/d}$, $T_{\min}=\sqrt{d(d-2)}/(2\pi R)$. 高于最低温度可有两种黑洞半径. 小半径处 $\dd T/\dd r_h<0$, 大半径处为正; 又 $\dd M/\dd r_h>0$, 所以小黑洞分支热容负、大黑洞分支热容正. 三个候选是三个鞍点, 不能把有负热容者也称为同样稳定的正则平衡相.

\originalsection{实际比较自由能而不是只数解}
相对同一热 AdS vacuum/Casimir 能量减去公共项, 黑洞质量与熵为
\[
M_{\rm rel}=\frac{(d-1)\Omega_{d-1}}{16\pi G_{d+1}}
r_h^{d-2}(1+r_h^2/R^2),\quad
S=\frac{\Omega_{d-1}r_h^{d-1}}{4G_{d+1}}.
\]
由 $F=M_{\rm rel}-TS$ 直接得
\[
F_{\rm BH}=\frac{\Omega_{d-1}}{16\pi G_{d+1}}
r_h^{d-2}(1-r_h^2/R^2).
\]
这相当于用合适的 boundary 反项或共同 vacuum subtraction 比较欧氏作用量, 不必把无限体积发散各自当成物理自由能. Classical 热 AdS 的相对自由能为0, 气体修正是 $O(1)$. 因而当 $r_h<R$, 黑洞 $F>0$, 热 AdS 优势; 当 $r_h>R$, 大黑洞的 $F<0$, 它优势.

两黑洞分支在最低温度处相遇, 此后大分支熵总更大, 所以 $\dd(F_{\rm big}-F_{\rm small})/\dd T=-(S_{\rm big}-S_{\rm small})<0$, 大分支自由能更低. 但在 $r_{h,\rm big}<R$ 区间, 即使它局部热容正, 仍输给热 AdS, 只是可能的亚稳鞍点. 全局优势要等到
\[
r_h=R,\qquad T_{\rm HP}=\frac{d-1}{2\pi R}
\]
才交换. 这是 Hawking--Page 相变. 自由能两支在交换处相等, \textbf{自由能本身连续}, 一阶导数 $-S$ 跳跃, 才称一阶相变. 热态的 $O(1)$ 与 $O(N^2)$ 部门区别不是说 $F$ 数值必须不连续.

\begin{center}
\begin{tikzpicture}
\begin{axis}[width=8cm,height=4.8cm,xlabel={$r_h/R$},ylabel={$TR$},
xmin=0.18,xmax=2,ymin=0.4,ymax=1,grid=major]
\addplot[structurecolor,thick] table[x=a,y=TR,col sep=comma]{qa/hawking-page.csv};
\addplot[only marks,mark=*,black] coordinates {(0.70710678,0.45015816)(1,0.47746483)};
\node[font=\small,anchor=north] at (axis cs:0.70710678,0.447) {$T_{\min}$};
\node[font=\small,anchor=south west] at (axis cs:1,0.47746483) {$T_{\rm HP}$};
\end{axis}
\end{tikzpicture}
\end{center}
图使用 $d=4$ 公式直接采样. 最低温度与自由能交换温度不同, 这也避免把“两解开始存在”误当成“相变发生”. 在有限边界体积, sharp 相变是大 $N$ 极限的热力学结构, 有限 $N$ 的解析性和修正需另外分析.

\originalsection{规范群不变单态投影是矩阵振子解释的关键}
考察两种频率为1的 adjoint 玻色振子. 若不作规范投影, 真是 $2N^2$ 个独立振子, 有
\[
\log Z_{\rm unprojected}=-2N^2\log(1-\ee^{-\beta}),
\]
任意固定有限 $\beta$ 都是 $O(N^2)$, \textbf{不能}据此推出低温自由能 $O(1)$. 想模拟规范 CFT 的物理态, 必须施加 Gauss 约束, 只留群不变单态.

在大 $N$、低能部门, 群不变单态由循环单迹字词及它们的多迹组合生成. 长度 $\ell$ 的两字母单迹数粗略是 $2^\ell/\ell$, 所以单迹热和包含 $\sum_\ell(2\ee^{-\beta})^\ell/\ell$. 当 $\beta>\log2$, 可收敛成不随 $N^2$ 增长的有限低温部门; 当 $\beta$ 接近 $\log2$, 指数态密度使这种低能展开失效.

群不变单态的定义是 $\widehat U(g)|\Psi\rangle=|\Psi\rangle$ 对全部群元素 $g$ 成立. 在这个矩阵量子力学模型, 对物理态施加 Gauss 约束等价于这份投影, 不是只保留每个矩阵某一条对角线.

更具体, 投影可用热 holonomy $U$ 的群积分实现. 单字母函数是 $z(q)=2q$, $q=\ee^{-\beta}$. 具有 holonomy 的两种玻色振子热迹给 $\log Z(U)=\sum_n(2q^n/n)|\operatorname{Tr}U^n|^2$. Haar 本征值测度含 $\prod_{i<j}|\ee^{\ii\theta_i}-\ee^{\ii\theta_j}|^2$, 对其对数用 Fourier 展开, 有 $S_{\rm Haar}=\sum_n|\operatorname{Tr}U^n|^2/n$ 加独立于分布的常数. 两项相减, 共同给有效作用量
\[
S_{\rm eff}[U]=\sum_{n\ge1}\frac{1-2q^n}{n}
|\operatorname{Tr}U^n|^2
\]
加上不影响主要鞍点的常数. 低温所有系数正, 均匀本征值分布 $\operatorname{Tr}U^n\simeq0$ 有优势; 第一系数在 $q=1/2$ 变号, 出现聚集本征值、$O(N^2)$ 部门的转变. 这展示为何规范 constraint 和巨大态密度共同决定热行为, 而不是把未投影振子总数直接当作低温物理态数.

真实球面 SYM 的模式有不同频率、相互作用和群不变单态约束, 转变温度也不由这个两字母模型精确预测. 它只给熵与能量竞争的机制: 低温群不变单态 $O(1)$ 部门对应热 AdS, 高能 $O(N^2)$ 部门对应黑洞. 强耦合具体 $T_{\rm HP}$ 则由上面实际引力自由能决定. 平直空间极限 $R_b\to\infty$ 在固定正 $T$ 下处于这一高温部门, 没有由有限球半径产生的同样非零温度阈值.

\originalsection{化学势是守恒荷的源}
规定正 $\mu$ 提高正荷态权重, 巨正则配分函数为
\[
\Xi(T,\mu)=\operatorname{Tr}\ee^{-\beta(H-\mu Q)},
\qquad \dd f=-s\dd T-\rho\dd\mu.
\]
若文献采用相反的化学势号, 需要同时改变其源定义. 对 $\mathcal N=4$ 的一个适当 $U(1)$ R 对称流, 源是边界 $A_t$. 正则欧氏视界的时间圆可收缩, 通常选规范令 $A_t(z_h)=0$, 化学势是
\[
\mu=A_t(0)-A_t(z_h).
\]
这个势差比未经指定规范的单点 $A_t$ 更有明确物理意义. 求解后电场 radial 通量给 charge 密度, 而不是把 $\mu$ 本身误称为电荷.

\originalsection{一个完整的 Einstein--Maxwell 例子}
为固定系数, 明确选择
\[
S=\frac1{16\pi G_{d+1}}\int\sqrt{-g}\dd^{d+1}X
\left[\mathcal R+\frac{d(d-1)}{R^2}-\frac{R^2}{4}F_{MN}F^{MN}\right].
\]
此处 $A$ 有 $L^{-1}$ 量纲, 规范项 $R^2F^2$ 与曲率同量纲. 这是一个确定的有效截断; 某些实际 IIB Cartan 选择还会源标量, 需按一致截断增补, 不由这一个模型覆盖全部 R-charge 相.

取前述 planar 假设形式, $A_t=A_t(z)$, $d>2$. Maxwell 方程给 $\partial_z(z^{3-d}A_t')=0$, 所以 $A_t=\mu-Q_0z^{d-2}$. Spatial Einstein 分量给
\[
zf'-df=-d+\frac{(d-2)^2Q_0^2}{2(d-1)}z^{2d-2}.
\]
积分得
\[
f(z)=1-M_0z^d+\frac{(d-2)Q_0^2}{2(d-1)}z^{2d-2}.
\]
$Q_0$ 是电场积分常数, 其物理密度还带作用量的归一化. 将 $d=4$ 且 $A_t(z_h)=0$ 代入, 定义 $q^2=\mu^2z_h^2/3$, 得
\[
f=1-(1+q^2)(z/z_h)^4+q^2(z/z_h)^6,
\qquad A_t=\mu[1-(z/z_h)^2],
\]
以及
\[
T=\frac{1-q^2/2}{\pi z_h},\qquad
\mu=\frac{\sqrt3q}{z_h},\qquad
s=\frac{R^3}{4G_5z_h^3}.
\]
$q^2<2$ 是这一非极端正温度分支, $q^2=2$ 的经典极端解仍有有限面积熵. 是否对全部额外带荷场稳定, 不是只看这份温度式就能断定.

设 $C=R^3/(16\pi G_5)$, 同一归一化下的密度是
\[
e=\frac{3C(1+q^2)}{z_h^4},\quad
p=\frac{C(1+q^2)}{z_h^4},\quad
\rho=\frac{2\sqrt3Cq}{z_h^3},\quad
s=\frac{4\pi C}{z_h^3}.
\]
直接微分 $q,z_h$ 可验证 $\dd e=T\dd s+\mu\dd\rho$, 也有 $e+p=Ts+\mu\rho$. 这同时检查温度、密度和作用量系数没有混用. 精确的有限密度物质相、输运及可能的标量凝聚仍需实际多场作用量和稳定性分析, 本讲完成的是该明确 EM 模型及源字典的推导.
